\documentclass[11pt,a4paper]{article}
% Wspólny szablon dla kart supermocy (jedna karta = jedna strona A4)
% Użycie w karcie: \documentclass[11pt,a4paper]{article} % Wspólny szablon dla kart supermocy (jedna karta = jedna strona A4)
% Użycie w karcie: \documentclass[11pt,a4paper]{article} % Wspólny szablon dla kart supermocy (jedna karta = jedna strona A4)
% Użycie w karcie: \documentclass[11pt,a4paper]{article} \input{../preambula} \begin{document}

% Silnik tectonic (XeTeX): polskie znaki przez fontspec, nie przez inputenc/T1
\usepackage{fontspec}
\setmainfont{Arial}
\usepackage[polish]{babel}
\usepackage[a4paper,margin=1.5cm]{geometry}
\usepackage{amsmath,amssymb}
\usepackage{tikz}
\usetikzlibrary{arrows.meta,positioning,calc,shapes.geometric}
\usepackage[most]{tcolorbox}
\usepackage{enumitem}
\usepackage{booktabs}
\usepackage{xcolor}

\pagestyle{empty}
\setlength{\parindent}{0pt}
\setlength{\parskip}{0.4em}

\definecolor{cel}{HTML}{1F4E79}
\definecolor{regula}{HTML}{E8F1FA}
\definecolor{pulapka}{HTML}{FDECEA}
\definecolor{okres}{HTML}{C0392B}
\definecolor{skonczony}{HTML}{1E8449}

% Ramka na regułę
\newtcolorbox{regula}[1][]{
  colback=regula, colframe=cel, boxrule=0.8pt, arc=2pt,
  left=6pt, right=6pt, top=4pt, bottom=4pt, title=#1,
  fonttitle=\bfseries, coltitle=white
}

% Ramka na pułapkę z kartkówki
\newtcolorbox{pulapka}[1][]{
  colback=pulapka, colframe=okres, boxrule=0.8pt, arc=2pt,
  left=6pt, right=6pt, top=4pt, bottom=4pt, title=#1,
  fonttitle=\bfseries, coltitle=white
}

% Nagłówek karty: numer, tytuł, cel
\newcommand{\naglowekKarty}[3]{%
  {\LARGE\bfseries\color{cel} Karta #1: #2}\\[2pt]
  {\large Cel: #3}\par
  \vspace{0.3em}\noindent\rule{\linewidth}{1pt}\par
}
 \begin{document}

% Silnik tectonic (XeTeX): polskie znaki przez fontspec, nie przez inputenc/T1
\usepackage{fontspec}
\setmainfont{Arial}
\usepackage[polish]{babel}
\usepackage[a4paper,margin=1.5cm]{geometry}
\usepackage{amsmath,amssymb}
\usepackage{tikz}
\usetikzlibrary{arrows.meta,positioning,calc,shapes.geometric}
\usepackage[most]{tcolorbox}
\usepackage{enumitem}
\usepackage{booktabs}
\usepackage{xcolor}

\pagestyle{empty}
\setlength{\parindent}{0pt}
\setlength{\parskip}{0.4em}

\definecolor{cel}{HTML}{1F4E79}
\definecolor{regula}{HTML}{E8F1FA}
\definecolor{pulapka}{HTML}{FDECEA}
\definecolor{okres}{HTML}{C0392B}
\definecolor{skonczony}{HTML}{1E8449}

% Ramka na regułę
\newtcolorbox{regula}[1][]{
  colback=regula, colframe=cel, boxrule=0.8pt, arc=2pt,
  left=6pt, right=6pt, top=4pt, bottom=4pt, title=#1,
  fonttitle=\bfseries, coltitle=white
}

% Ramka na pułapkę z kartkówki
\newtcolorbox{pulapka}[1][]{
  colback=pulapka, colframe=okres, boxrule=0.8pt, arc=2pt,
  left=6pt, right=6pt, top=4pt, bottom=4pt, title=#1,
  fonttitle=\bfseries, coltitle=white
}

% Nagłówek karty: numer, tytuł, cel
\newcommand{\naglowekKarty}[3]{%
  {\LARGE\bfseries\color{cel} Karta #1: #2}\\[2pt]
  {\large Cel: #3}\par
  \vspace{0.3em}\noindent\rule{\linewidth}{1pt}\par
}
 \begin{document}

% Silnik tectonic (XeTeX): polskie znaki przez fontspec, nie przez inputenc/T1
\usepackage{fontspec}
\setmainfont{Arial}
\usepackage[polish]{babel}
\usepackage[a4paper,margin=1.5cm]{geometry}
\usepackage{amsmath,amssymb}
\usepackage{tikz}
\usetikzlibrary{arrows.meta,positioning,calc,shapes.geometric}
\usepackage[most]{tcolorbox}
\usepackage{enumitem}
\usepackage{booktabs}
\usepackage{xcolor}

\pagestyle{empty}
\setlength{\parindent}{0pt}
\setlength{\parskip}{0.4em}

\definecolor{cel}{HTML}{1F4E79}
\definecolor{regula}{HTML}{E8F1FA}
\definecolor{pulapka}{HTML}{FDECEA}
\definecolor{okres}{HTML}{C0392B}
\definecolor{skonczony}{HTML}{1E8449}

% Ramka na regułę
\newtcolorbox{regula}[1][]{
  colback=regula, colframe=cel, boxrule=0.8pt, arc=2pt,
  left=6pt, right=6pt, top=4pt, bottom=4pt, title=#1,
  fonttitle=\bfseries, coltitle=white
}

% Ramka na pułapkę z kartkówki
\newtcolorbox{pulapka}[1][]{
  colback=pulapka, colframe=okres, boxrule=0.8pt, arc=2pt,
  left=6pt, right=6pt, top=4pt, bottom=4pt, title=#1,
  fonttitle=\bfseries, coltitle=white
}

% Nagłówek karty: numer, tytuł, cel
\newcommand{\naglowekKarty}[3]{%
  {\LARGE\bfseries\color{cel} Karta #1: #2}\\[2pt]
  {\large Cel: #3}\par
  \vspace{0.3em}\noindent\rule{\linewidth}{1pt}\par
}


\begin{document}

\naglowekKarty{1}{Minuty na godziny w ułamku dziesiętnym}{zamieniać minuty na ułamek godziny bez mylenia z setnymi}

\vspace{0.3em}
\begin{regula}[REGUŁA]
Godzina ma $60$ minut, \emph{nie} $100$. Dlatego
$$t\ \text{min} = \frac{t}{60}\ \text{h}$$
\textbf{Najważniejsze:} $15\ \text{min}\neq 0{,}15\ \text{h}$, bo $\dfrac{15}{60}=\dfrac{1}{4}=0{,}25$.
\end{regula}

\vspace{0.4em}
\textbf{Druga droga: dzielenie pod kreską.} Zamiast pamiętać wynik, podziel minuty przez $60$:
$$18 : 60 = 0{,}3 \qquad\text{(bo } 180 : 60 = 3\text{, czyli } 18 \text{ dziesiątych)}$$
$$45 : 60 = 0{,}75 \qquad\text{(bo } 450 : 60 = 7 \text{ r. } 30,\ 300 : 60 = 5\text{)}$$

\begin{center}
\begin{tikzpicture}[>=Stealth, x=0.2cm]
  % Skala minut (góra) i skala godzin w ułamku dziesiętnym (dół)
  \draw[->] (-0.5,0) -- (62,0) node[right, font=\small] {min};
  \draw[->] (-0.5,-2.6) -- (62,-2.6) node[right, font=\small] {h};
  \foreach \m in {0,6,...,60} {
    \draw (\m,0.15) -- (\m,-0.15);
    \node[above, font=\scriptsize] at (\m,0.15) {\m};
  }
  \foreach \k in {0,1,...,10} {
    \pgfmathsetmacro{\xx}{6*\k}
    \draw (\xx,-2.45) -- (\xx,-2.75);
    \pgfmathsetmacro{\lab}{0.1*\k}
    \node[below, font=\scriptsize] at (\xx,-2.7) {\pgfmathprintnumber[fixed,precision=1]{\lab}};
  }
  % Zaznaczenie 18 min = 0,3 h oraz 45 min = 0,75 h
  \fill[cel] (18,0) circle (2.2pt);
  \fill[cel] (45,0) circle (2.2pt);
  \draw[dashed, cel] (18,0) -- (18,-2.6);
  \draw[dashed, cel] (45,0) -- (45,-2.6);
  \node[font=\small\bfseries, color=cel] at (18,-3.5) {$18\ \text{min}=0{,}3\ \text{h}$};
  \node[font=\small\bfseries, color=cel] at (45,-3.5) {$45\ \text{min}=0{,}75\ \text{h}$};
  \node[font=\small, anchor=west] at (-0.5,1.1) {\textbf{Klocek:} $6\ \text{min}=0{,}1\ \text{h}$ (jedna dziesiąta godziny)};
\end{tikzpicture}
\end{center}

\vspace{0.2em}
\textbf{Tabela klocków (zapamiętaj 6 i 15):}
\begin{center}
\renewcommand{\arraystretch}{1.9}
\begin{tabular}{@{}lll@{}}
\toprule
Minuty & Ułamek zwykły & Dziesiętny \\
\midrule
$6$ & $\dfrac{6}{60}=\dfrac{1}{10}$ & $0{,}1$ \\
$15$ & $\dfrac{15}{60}=\dfrac{1}{4}$ & $0{,}25$ \\
$30$ & $\dfrac{30}{60}=\dfrac{1}{2}$ & $0{,}5$ \\
$45$ & $\dfrac{45}{60}=\dfrac{3}{4}$ & $0{,}75$ \\
\bottomrule
\end{tabular}
\end{center}

\vspace{0.2em}
\textbf{Przykłady rozstrzygnięte:}
\begin{itemize}[leftmargin=*, itemsep=1pt]
  \item $18\ \text{min} = 3 \cdot 6\ \text{min} = 3\cdot 0{,}1 = \mathbf{0{,}3\ h}$
  \item $2\ \text{godz}\ 15\ \text{min} = 2 + 0{,}25 = \mathbf{2{,}25\ h}$
  \item $1\ \text{godz}\ 12\ \text{min} = 1 + 2\cdot 0{,}1 = \mathbf{1{,}2\ h}$
  \item $27\ \text{min} = 24 + 3 = 4\cdot 0{,}1 + 0{,}05 = \mathbf{0{,}45\ h}$
\end{itemize}

\begin{pulapka}[PUŁAPKI Z KARTKÓWKI]
\begin{enumerate}[leftmargin=*, itemsep=2pt]
  \item $15\ \text{min}\neq 0{,}15\ \text{h}$. Kwadrans to $0{,}25\ \text{h}$.
  \item $2\ \text{godz}\ 15\ \text{min}\neq 2{,}15\ \text{h}$. Zapis $2{,}15$ to $2$ i $15$ setnych, a $15$ minut to $\dfrac{15}{60}=0{,}25$.
  \item Nie dziel przez $100$. Godzina to $60$ minut.
\end{enumerate}
\end{pulapka}

\end{document}
